\clearpage\thispagestyle{empty}
\pdfbookmark[0]{Sources for Instruments (20 November 1820)}{instrumentsnov1820-sources}
\begin{center}\textbf{Sources for Instruments (20 November 1820)}\end{center}
\noindent This AI translation follows \textit{Instrumente} in \href{https://gdz.sub.uni-goettingen.de/id/PPN236007467}{\textit{Werke}, vol.~VI} (Göttingen, 1874), p.~435. The first printing appeared under \textit{Göttingen} in the \textit{Göttingische gelehrte Anzeigen}, no.~187, 20 November 1820, \href{https://images.p.sub.uni-goettingen.de/iiif/image/gdz:PPN319721507_1820:00001879/full/1600,/0/default.jpg}{p.~1865} and \href{https://images.p.sub.uni-goettingen.de/iiif/image/gdz:PPN319721507_1820:00001880/full/1600,/0/default.jpg}{p.~1866}, in the Göttingen State and University Library copy.
\par\smallskip\noindent The source calls Hardy’s clock a \textit{Tertienuhr}. A \href{https://www.duden.de/rechtschreibung/Tertie}{\textit{Tertie}} is one sixtieth of a second. The words about the pendulum’s sensitivity remain legible beneath the stray strokes in the collected edition and are also present in the first printing. Translation and typesetting produced by AI.
\clearpage
